\section{稠密子集上的映射延拓}
\label{sec:mmap-extension}

只在稠密子集上知道一个映射的值，是常见的情形：先在有理数上定义一个运算，再希望把它推广到全体实数；或者只在有限精度的网格上做过测量，希望确定一个连续量。设$D\subseteq X$稠密，$f:D\to Y$已知。对$x\in X\setminus D$，自然的办法是取$D$中趋于$x$的点列$d_n$，把$f(d_n)$的极限定义为新值。这要求两件事：像点列确实有极限，并且换一个逼近列不会得到不同的结果。本节说明，目标空间完备保证前者，均匀连续保证后者，并且延拓后保留原来的误差控制。

\subsection{延拓的构造}
\label{sec:mmap-extension-construction}

我们先给出定理，再逐步说明每个条件在构造中的作用。

\begin{theorem}[均匀连续映射的延拓]
\label{thm:mmap-extension}
设$D$是度量空间$X$的稠密子集，$Y$是完备度量空间，$f:D\to Y$均匀连续。则存在唯一的连续映射$F:X\to Y$，使$F|_D=f$。
\end{theorem}

\begin{proof}
定义。对$x\in X$，由稠密性和第~\ref{chap:metric-spaces}章的闭包序列判据，取$d_n\in D$使$d_n\to x$。收敛点列是柯西列，均匀连续映射把柯西列送到柯西列，所以$(f(d_n))$是$Y$中的柯西列，由$Y$完备收敛。令$F(x)=\lim_nf(d_n)$。

与逼近列无关。设另有$d_n'\in D$趋于$x$。把两列交替排成新点列$d_1,d_1',d_2,d_2',\ldots$，它仍趋于$x$，于是它的像也是柯西列并收敛到某个$z$。$(f(d_n))$与$(f(d_n'))$都是这个像点列的子列，因而都趋于$z$。所以$F(x)$不依赖于逼近列的选择。

延拓。若$x\in D$，可以取常值逼近列$d_n=x$，得$F(x)=f(x)$。

唯一性。设$G:X\to Y$连续且$G|_D=f$。对$x\in X$及趋于$x$的$d_n\in D$，由序列判据$G(x)=\lim G(d_n)=\lim f(d_n)=F(x)$。$F$本身的连续性在下一小节与误差控制一起证明。
\end{proof}

唯一性的论证只用到$D$稠密与$G$连续，不需要$Y$完备或$f$均匀连续：两个连续映射若在稠密子集上相同，就处处相同。存在性则两个条件都要用，缺一个都会失败。

缺少目标空间的完备性时，像点列可能没有极限。取$X=\mathbb R$、$D=\mathbb Q$、$Y=\mathbb Q$，$f$为$\mathbb Q$上的恒等映射，它是等距映射，当然均匀连续。若有连续延拓$F:\mathbb R\to\mathbb Q$，沿$\sqrt2$的有理截断$q_n$应有$F(\sqrt2)=\lim q_n=\sqrt2$，而$\sqrt2\notin\mathbb Q$。完备性要求作用在目标空间上，定义域$X$完备与否无关紧要。

缺少均匀连续时，不同的逼近列可能给出不同的极限。在$D=\mathbb Q\cap[0,1]$上令$f(q)=0$（$q<1/\sqrt2$）、$f(q)=1$（$q>1/\sqrt2$）。$1/\sqrt2$不是有理数，所以$f$在$D$的每一点附近都是常值，因而连续；但从左、右两侧逼近$1/\sqrt2$的有理数列，像分别恒为$0$与$1$，没有连续延拓到$[0,1]$。这正说明了$f$不均匀连续：两侧的点可以任意接近，函数值却始终相差$1$。

\subsection{延拓后的误差控制}
\label{sec:mmap-extension-error}

延拓不仅存在，还继承原映射在$D$上的误差控制，阈值与常数都不变坏。

\begin{theorem}[延拓保持均匀连续与Lipschitz常数]
\label{thm:mmap-extension-control}
在定理~\ref{thm:mmap-extension}的条件下，设$\varepsilon,\delta>0$满足：$d_1,d_2\in D$且$d(d_1,d_2)<\delta$时$d(f(d_1),f(d_2))<\varepsilon$。则$x,x'\in X$且$d(x,x')<\delta$时$d(F(x),F(x'))\leq\varepsilon$。特别地，$F$均匀连续；若$f$以$L$为Lipschitz常数，则$F$也以$L$为Lipschitz常数。
\end{theorem}

\begin{proof}
取$d_n\to x$，$d_n'\to x'$，$d_n,d_n'\in D$。由三角不等式，$d(d_n,d_n')\to d(x,x')<\delta$，所以$n$充分大时$d(d_n,d_n')<\delta$，从而$d(f(d_n),f(d_n'))<\varepsilon$。令$n\to\infty$，由距离的连续性得$d(F(x),F(x'))\leq\varepsilon$。对任意$\varepsilon$，用$\varepsilon/2$对应的$\delta$，就得到$F$的均匀连续性，特别地$F$连续。

若对$D$中所有点都有$d(f(d_1),f(d_2))\leq L\,d(d_1,d_2)$，则对每个$n$有$d(f(d_n),f(d_n'))\leq L\,d(d_n,d_n')$，两边取极限得$d(F(x),F(x'))\leq L\,d(x,x')$。
\end{proof}

由这两个定理，哪些函数可以由稠密观测唯一确定，答案很清楚。在有理数上给定的函数若在每个有界区间上均匀连续，就唯一地确定了实轴上的连续函数：在每个$[-M,M]$上分别延拓，由唯一性，不同区间上的延拓在公共部分一致；例如多项式$p(q)$在$\mathbb Q\cap[-M,M]$上Lipschitz连续，延拓结果就是同一个多项式在实轴上的取值。反过来，若只知道目标函数连续，在有理点上的值就足以确定它；若函数可以不连续，有理点上的值什么也确定不了，例如在有理点取$0$、在无理点任意取值的函数。

延拓定理还完成了第~\ref{chap:metric-spaces}章留下的一步：完备化在等距意义下唯一。

\begin{theorem}[完备化的唯一性]
\label{thm:mmap-completion-unique}
设$\iota:X\to\hat X$与$j:X\to\tilde X$都是保持距离的映射，$\hat X,\tilde X$完备，像$\iota(X)$与$j(X)$分别稠密。则存在满射等距映射$\Phi:\hat X\to\tilde X$，使$\Phi\circ\iota=j$。
\end{theorem}

\begin{proof}
在稠密子集$\iota(X)$上定义$\phi(\iota(x))=j(x)$，它保持距离，因而是Lipschitz常数为$1$的均匀连续映射。$\tilde X$完备，由定理~\ref{thm:mmap-extension}，$\phi$延拓为连续映射$\Phi:\hat X\to\tilde X$。对$\xi,\eta\in\hat X$，取$\iota(x_n)\to\xi$，$\iota(y_n)\to\eta$，则
\begin{equation*}
\begin{aligned}
d(\Phi(\xi),\Phi(\eta))&=\lim d(j(x_n),j(y_n))=\lim d(x_n,y_n)\\
&=\lim d(\iota(x_n),\iota(y_n))=d(\xi,\eta),
\end{aligned}
\end{equation*}
所以$\Phi$保持距离。像$\Phi(\hat X)$与完备空间$\hat X$等距，因而完备，在$\tilde X$中是闭集；它又包含稠密集$j(X)$，所以等于$\tilde X$。
\end{proof}
